\documentclass[11pt]{article}

\usepackage[margin=1in]{geometry}
\usepackage[T1]{fontenc}
\usepackage[utf8]{inputenc}
\usepackage{hyperref}
\usepackage{longtable}
\usepackage{array}

\hypersetup{
  colorlinks=true,
  linkcolor=blue,
  urlcolor=blue
}

\newcommand{\code}[1]{\texttt{#1}}

\title{Datamak Workflow: From Equilibrium Ingestion to Surrogate Training}
\author{}
\date{Updated: 2026-03-26}

\begin{document}

\sloppy
\maketitle
\tableofcontents

\section{Purpose}

This document describes the main Datamak workflow from equilibrium ingestion to
surrogate-model training. It is intended as a bridge between the high-level
architecture notes and the concrete scripts used in daily operation.

The scope is the path
\[
\code{data\_equil} \rightarrow \code{gk\_study} \rightarrow \code{gk\_input}
\rightarrow \code{gk\_run} \rightarrow \code{gk\_surrogate}
\]
with emphasis on the repository entry points that implement each step.

\section{Central Database and Core Tables}

The main SQLite database is \code{gyrokinetic\_simulations.db} at repository
root. The baseline schema is created by
\code{database/create\_gyrokinetic\_db.py}.

\subsection{Key tables}

\begin{longtable}{|>{\raggedright\arraybackslash}p{0.22\textwidth}|>{\raggedright\arraybackslash}p{0.70\textwidth}|}
\hline
\textbf{Table} & \textbf{Role} \\
\hline
\code{data\_origin} & Metadata describing where equilibria came from. \\
\hline
\code{data\_equil} & Raw equilibrium/profile inputs, including stored \code{pfile}, \code{gfile}, or \code{transpfile} references and content. \\
\hline
\code{gk\_code} & Gyrokinetic code identity such as GX or CGYRO. \\
\hline
\code{gk\_model} & Active input templates and model flags used to generate GK inputs. \\
\hline
\code{gk\_study} & A study linking one equilibrium to one gyrokinetic code. \\
\hline
\code{gk\_input} & Generated gyrokinetic inputs, their content, radial coordinate, parsed physics columns, and workflow status. \\
\hline
\code{gk\_batch} & Metadata for exported remote batch databases. \\
\hline
\code{gk\_run} & Remote execution records and output metrics such as \code{gamma\_max}, \code{ky\_abs\_mean}, and \code{diffusion}. \\
\hline
\code{gk\_surrogate} & Registry of trained surrogate artifacts and metadata. \\
\hline
\code{sg\_estimate} & Surrogate predictions attached to existing \code{gk\_input} rows. \\
\hline
\end{longtable}

\subsection{Operational status contracts}

Datamak uses status fields as workflow contracts between CLI scripts, the GUI,
and remote execution.

\begin{itemize}
\item \code{gk\_input.status}: typically \code{NEW}, \code{WAIT},
\code{TORUN}, \code{BATCH}, \code{CRASHED}, \code{SUCCESS}, \code{ERROR}.
\item \code{gk\_batch.status}: typically \code{CREATED}, \code{SENT},
\code{LAUNCHED}, \code{SYNCED}.
\item \code{gk\_run.status}: workflow-dependent, commonly \code{TORUN},
\code{RUNNING}, \code{SUCCESS}, \code{CONVERGED}, \code{CRASHED}.
\end{itemize}

\section{Workflow Overview}

\subsection{Step 0: Bootstrap the schema}

The repository bootstrap script
\code{database/create\_gyrokinetic\_db.py} creates the main schema and applies
guarded additive upgrades, including:

\begin{itemize}
\item \code{data\_origin} and \code{data\_equil} for source data,
\item \code{gk\_study}, \code{gk\_model}, and \code{gk\_input} for GK setup,
\item \code{gk\_batch} and \code{gk\_run} for execution tracking,
\item supporting indexes and compatibility \code{ALTER TABLE} logic.
\end{itemize}

Typical command:

\begin{verbatim}
python3 database/create_gyrokinetic_db.py --db gyrokinetic_simulations.db
\end{verbatim}

\subsection{Step 1: Ingest equilibria into \code{data\_equil}}

There are two main ingestion families.

\subsubsection{Mate Kinetic EFIT path}

The script \code{db\_update/populate\_data\_equil\_from\_Mate\_KinEFIT.py}:

\begin{enumerate}
\item scans a root directory for subfolders,
\item finds \code{p*} and \code{g*} files,
\item reads file contents into memory,
\item creates or reuses a \code{data\_origin} row,
\item inserts new \code{data\_equil} rows,
\item runs \code{update\_activate\_Mate\_KinEFIT.sql} to activate selected rows.
\end{enumerate}

In this path, the database stores both filenames and raw text content of the
\code{pfile} and \code{gfile}. This makes later GK input generation
self-contained.

\subsubsection{TRANSP path}

The TRANSP-based workflows populate \code{data\_equil} using CDF metadata
instead of local \code{pfile}/\code{gfile} pairs.

\begin{itemize}
\item Semi-automatic path: \code{db\_update/populate\_data\_equil\_from\_Alexei\_Transp\_09.py}
and related scripts.
\item Flux full-auto path: \code{db\_update/Transp\_full\_auto/MainSteps\_1\_launch\_on\_laptop.sh},
\code{MainSteps\_2\_launch\_on\_flux.sh}, and
\code{MainSteps\_3\_launch\_on\_laptop.sh}.
\end{itemize}

The Flux full-auto path uses a temporary Flux-local database, populates
\code{data\_equil} and \code{transp\_timeseries} there, generates inputs
remotely, then syncs the results back into the main database.

\subsection{Step 2: Create studies that connect equilibria to a GK code}

\code{gk\_study} is the bridge between one equilibrium and one gyrokinetic
code. The later input-generation scripts assume that:

\begin{itemize}
\item the relevant \code{data\_equil} rows are active,
\item the appropriate \code{gk\_code} row exists,
\item one or more \code{gk\_model} rows are marked active.
\end{itemize}

In practice, studies may be created by SQL activation scripts or by the
full-auto TRANSP workflow.

\subsection{Step 3: Generate \code{gk\_input} rows and input content}

The next stage converts equilibrium data plus template settings into explicit
gyrokinetic inputs.

\subsubsection{Pfile/Gfile-driven input generation}

The script
\code{db\_update/create\_gk\_input\_from\_pyrokinetic\_with\_pfile\_and\_gfile.py}
operates on active Mate equilibria. Its main responsibilities are:

\begin{enumerate}
\item fetch active \code{gk\_study} rows and their stored \code{pfile}/\code{gfile} content,
\item fetch active \code{gk\_model} template definitions,
\item sweep over a \code{psin} grid,
\item call Pyrokinetics with the requested template and local-geometry settings,
\item insert or update \code{gk\_input} rows with generated input content.
\end{enumerate}

The inserted rows include:

\begin{itemize}
\item \code{gk\_study\_id},
\item \code{gk\_model\_id},
\item filename and output path,
\item full input \code{content},
\item \code{psin},
\item workflow \code{status}, usually starting at \code{WAIT}.
\end{itemize}

\subsubsection{TRANSP-driven input generation}

The script
\code{db\_update/create\_gk\_input\_from\_pyrokinetic\_with\_transpfile.py}
plays the same role for TRANSP data. Instead of local \code{pfile}/\code{gfile}
pairs, it loads the TRANSP source, uses Pyrokinetics to evaluate a requested
time slice, and writes the resulting GX-style input content into
\code{gk\_input.content}.

The Flux full-auto workflow performs this generation remotely, then syncs the
resulting rows back into the main database.

\subsection{Step 4: Backfill parsed physics columns from \code{gk\_input.content}}

After input content exists, Datamak extracts structured scalar features from
that text using \code{db\_update/backfill\_gk\_input\_physics.py}.

This script parses the following sections from each input:

\begin{itemize}
\item \code{[geometry]} for quantities such as \code{Rmaj}, \code{qinp},
\code{shat}, \code{akappa}, \code{tri}, and \code{betaprim},
\item \code{[physics]} for scalar physics settings,
\item \code{[species]} for electron and main-ion density, temperature, mass,
charge, gradients, and collisionality.
\end{itemize}

The extracted values populate columns on \code{gk\_input}. This step is
important because the surrogate workflow reads engineered features from columns,
not by reparsing raw input text at training time.

\subsection{Step 5: Promote selected inputs to execution}

Generated inputs do not directly become surrogate-training data. They first
need simulation output.

The usual progression is:

\begin{enumerate}
\item create \code{gk\_input} rows with status \code{WAIT},
\item promote selected rows to \code{TORUN},
\item export those rows to a batch database,
\item run them remotely,
\item sync \code{gk\_run} results back into the main database.
\end{enumerate}

\subsection{Step 6: Batch export and remote execution}

The execution workflow is centered on the \code{batch/} folder.

\subsubsection{Create a batch database}

\code{batch/create\_batch\_database.py --copy-torun}:

\begin{itemize}
\item selects \code{gk\_input.status = TORUN},
\item copies those cases into a timestamped batch database under
\code{batch/new/},
\item creates \code{gk\_run} rows in that batch database,
\item updates source \code{gk\_input.status} to \code{BATCH},
\item inserts a \code{gk\_batch} row with status \code{CREATED}.
\end{itemize}

\subsubsection{Deploy and execute remotely}

\code{batch/deploy\_batch\_large.py} transfers non-empty batch databases and
HPC helper scripts to Perlmutter, updates \code{gk\_batch} status through
\code{SENT} and \code{LAUNCHED}, and submits remote jobs.

On the remote side, the job scripts consume \code{gk\_run} rows, write input
files, execute the solver, and update \code{gk\_run.status} together with
quantities such as \code{gamma\_max}, \code{ky\_abs\_mean}, and
\code{diffusion}.

\subsubsection{Sync back to the main database}

\code{batch/check\_launched\_batches.py} and
\code{batch/monitor\_remote\_runs.py} synchronize the remote execution results
and produce health reports. At this point, the main database contains the
joined information needed for surrogate training:

\begin{itemize}
\item input-side features on \code{gk\_input},
\item target outputs on \code{gk\_run}.
\end{itemize}

\section{Surrogate Training Workflow}

\subsection{Training data selection}

The surrogate training script is \code{db\_surrogate/train\_gamma\_surrogate.py}.
It loads rows from a join of \code{gk\_run} and \code{gk\_input}.

By default it keeps only:

\begin{itemize}
\item rows whose \code{gk\_run.status} is \code{SUCCESS} or \code{CONVERGED},
\item rows with finite target values,
\item rows where all required input features can be assembled.
\end{itemize}

Optional filters exist for \code{data\_origin} name or identifier.

\subsection{Feature construction}

The current feature vector is:

\begin{verbatim}
temp_ratio, dens_ratio, electron_tprim, electron_fprim, ion_tprim, ion_fprim,
ion_vnewk, electron_vnewk, mass_ratio, Rmaj, qinp, shat, shift, akappa,
akappri, tri, tripri, betaprim
\end{verbatim}

The three ratio features are computed at training time from the stored
\code{gk\_input} columns:

\begin{itemize}
\item \code{temp\_ratio = ion\_temp / electron\_temp},
\item \code{dens\_ratio = ion\_dens / electron\_dens},
\item \code{mass\_ratio = ion\_mass / (electron\_mass * 1836)}.
\end{itemize}

\subsection{Target definition}

The default regression target is \code{gk\_run.gamma\_max}. The script can also
train against other \code{gk\_run.*} quantities through the \code{--mapsto}
argument, for example \code{gk\_run.ky\_abs\_mean} or \code{gk\_run.diffusion}.

\subsection{Model fitting and artifacts}

The current baseline model is a
\code{sklearn.ensemble.RandomForestRegressor}. The script supports:

\begin{itemize}
\item optional train/test splitting,
\item configurable tree count and leaf size,
\item optional \code{log1p} transform of the target.
\end{itemize}

Outputs are written to \code{db\_surrogate/models/} by default:

\begin{itemize}
\item a pickle payload containing the fitted model and metadata,
\item a JSON sidecar with non-binary metadata such as feature ranges, metrics,
status filter, and target mapping.
\end{itemize}

Representative command:

\begin{verbatim}
python3 db_surrogate/train_gamma_surrogate.py \
  --db gyrokinetic_simulations.db \
  --name gamma_v1
\end{verbatim}

\subsection{Registering and using a trained surrogate}

Model files alone are sufficient for offline prediction, but Datamak also
supports database-backed registration through \code{gk\_surrogate}. The GUI code
contains helper logic to create that table if needed and record metadata for a
trained model.

Once a \code{gk\_surrogate} row exists, the script
\code{db\_surrogate/estimate\_gamma\_surrogate.py} can:

\begin{enumerate}
\item load the registered model,
\item reconstruct the same feature vector for all complete \code{gk\_input} rows,
\item predict mean values and an uncertainty proxy,
\item upsert the results into \code{sg\_estimate}.
\end{enumerate}

\section{End-to-End Summary}

The complete operational chain is:

\begin{enumerate}
\item create or refresh the schema,
\item ingest equilibria into \code{data\_equil},
\item create active studies and templates,
\item generate \code{gk\_input.content} with Pyrokinetics,
\item backfill structured plasma and geometry columns onto \code{gk\_input},
\item run selected inputs remotely and sync \code{gk\_run} metrics,
\item train a surrogate from the joined \code{gk\_input} and \code{gk\_run} data,
\item optionally register the model in \code{gk\_surrogate} and populate
\code{sg\_estimate}.
\end{enumerate}

\section{Practical Command Reference}

\begin{verbatim}
python3 database/create_gyrokinetic_db.py --db gyrokinetic_simulations.db

python3 db_update/populate_data_equil_from_Mate_KinEFIT.py --db gyrokinetic_simulations.db

python3 db_update/create_gk_input_from_pyrokinetic_with_pfile_and_gfile.py \
  --db gyrokinetic_simulations.db

python3 db_update/backfill_gk_input_physics.py --db gyrokinetic_simulations.db

python3 batch/create_batch_database.py --copy-torun

python3 batch/deploy_batch_large.py

python3 db_surrogate/train_gamma_surrogate.py --db gyrokinetic_simulations.db --name gamma_v1

python3 db_surrogate/estimate_gamma_surrogate.py --db gyrokinetic_simulations.db --surrogate-id 1
\end{verbatim}

\section{Current Limitations}

\begin{itemize}
\item The base schema bootstrap does not yet create \code{gk\_surrogate}; that
table is still created lazily in the GUI path.
\item The exact entry point for \code{gk\_study} creation depends on which
ingestion branch is used.
\item Some Flux full-auto scripts still carry environment assumptions that
should be documented and validated more explicitly.
\end{itemize}

\end{document}
